popatrz na to, jak to jest zbudowane w latex, dodaj na tej samej zasadzie nasze pytania, niech to będą wszyskie, powklejaj tam odpowiedzi w Standardowy RAG, GraphRag (hybrid), i naszą agentową stukruę, daj blue print pod wnioski i obserwacje jezeli jestes w stanie cos wymyslic, staraj sie jak najbardziej profejonalenie \section*{Porównanie: Termin 1 cyklu rekrutacji na I stopień studiów stacjonarnych}

\begin{tcolorbox}[colback=blue!5!white,colframe=blue!75!black,title=Pytanie użytkownika]
    Jaki jest termin rozpoczęcia i zakończenia 1 cyklu rekrutacji na I stopień studiów stacjonarnych?
\end{tcolorbox}

\noindent
\begin{tabularx}{\textwidth}{@{} XX @{}}
    \toprule
    \textbf{Standardowy RAG} & \textbf{GraphRAG (Hybrid)} \\
    \midrule
    \small 
    Termin rozpoczęcia 1 cyklu rekrutacji na I stopień studiów stacjonarnych to 1.06, a zakończenia 23.07 do godz. 14:00.
    & 
    \small 
    Termin rozpoczęcia 1 cyklu rekrutacji na I stopień studiów stacjonarnych to 1 czerwca, a zakończenia 11 lipca do godz. 14:00. \\
    \bottomrule
\end{tabularx}

\begin{tcolorbox}[colback=green!5!white,colframe=green!50!black,title=Wnioski i obserwacje]
    \begin{itemize}
        \item \textbf{Wysoka precyzja GraphRAG:} System oparty na grafach bezbłędnie poradził sobie z uszczegółowionym pytaniem. Dzięki strukturze grafu (tabela $\rightarrow$ kontekst kierunku $\rightarrow$ cykl 1) zlokalizował dokładną datę z właściwej kolumny (1.06 - 11.07), w pełni poprawnie interpretując treść z obrazka.
        \item \textbf{Błąd dopasowania kontekstu (Standard RAG):} Standardowy RAG udzielił błędnej odpowiedzi. Silnik wektorowy odnalazł fragment innego dokumentu, który pasował semantycznie ("Cykl rekrutacyjny 1", "Złożenie i opłacenie"), ale dotyczył innej formy studiów (gdzie data to 23.07). Wynika to z faktu, że wektory słabo radzą sobie z rygorystycznym wiązaniem wielu warunków naraz (I stopień + stacjonarne + cykl 1).
        \item \textbf{Interpretacja pojęcia "zakończenia":} Oba modele domyślnie zinterpretowały "zakończenie rekrutacji" jako zamknięcie systemu e-Rekrutacja, a nie ostateczny termin wpisu na listę studentów (kończący się 18.07).
    \end{itemize}
\end{tcolorbox}
\vspace{1em}

\newpage

\section*{Porównanie: Punkty za osiągnięcia sportowe w rekrutacji}

\begin{tcolorbox}[colback=blue!5!white,colframe=blue!75!black,title=Pytanie użytkownika]
    Za jakie osiągnięcia sportowe mogę zdobyć dodatkowe punkty aplikując na studia i ile tych punktów maksymalnie mogę dostać?
\end{tcolorbox}

\noindent
\begin{tabularx}{\textwidth}{@{} XX @{}}
    \toprule
    \textbf{Standardowy RAG} & \textbf{GraphRAG (Hybrid)} \\
    \midrule
    \small 
    Możesz zdobyć dodatkowe punkty za osiągnięcia sportowe, a maksymalnie możesz otrzymać do 150 punktów.
    & 
    \small 
    Możesz zdobyć dodatkowe punkty aplikując na studia za osiągnięcia sportowe w różnych dyscyplinach, takich jak strzelectwo, żeglarstwo, wspinaczka, wioślarstwo, unihokej, trójbój siłowy, tenis, szermierka, szachy, i wiele innych. Maksymalnie możesz zdobyć 300 punktów za każde z tych osiągnięć, w zależności od posiadanej klasy sportowej:
    \begin{itemize}
        \item 300 punktów - klasa mistrzowska międzynarodowa lub klasa mistrzowska
        \item 150 punktów - pierwsza klasa sportowa
        \item 100 punktów - druga klasa sportowa
    \end{itemize}
    Dodatkowo, za osiągnięcia w sportach olimpijskich, paralimpijskich lub w Akademickich Mistrzostwach Polski, punkty przyznawane są w zależności od miejsca zajętego w zawodach. \\
    \bottomrule
\end{tabularx}

\begin{tcolorbox}[colback=green!5!white,colframe=green!50!black,title=Wnioski i obserwacje]
    \begin{itemize}
        \item \textbf{Synergia danych w GraphRAG (Multi-document reasoning):} Graf wiedzy poprawnie połączył informacje. Model wymienił konkretne dyscypliny i podał prawidłową wartość maksymalną (300 punktów).
        \item \textbf{Uporczywy błąd standardowego RAG:} Zwykły RAG popełnił błąd. Silnik wektorowy "łapie" fragment z wyciągniętym ciągiem znaków "150 pkt" (dotyczącym innej tabelki), ignorując zupełnie właściwe zasady, co udowadnia, że RAG oparty tylko na wektorach jest bardzo podatny na zmyłki semantyczne.
    \end{itemize}
\end{tcolorbox}
\vspace{1em}

\newpage

\section*{Porównanie: Wymagane przedmioty maturalne (Automatyka i Robotyka)}

\begin{tcolorbox}[colback=blue!5!white,colframe=blue!75!black,title=Pytanie użytkownika]
    Jakie przedmioty maturalne są wymagane na kierunku Automatyka i Robotyka?
\end{tcolorbox}

\noindent
\begin{tabularx}{\textwidth}{@{} XX @{}}
    \toprule
    \textbf{Standardowy RAG} & \textbf{GraphRAG (Hybrid)} \\
    \midrule
    \small 
    Aby dostać się na Automatykę i Robotykę, potrzebne są przedmioty maturalne: matematyka albo fizyka albo informatyka. Dodatkowo, w drugim przedmiocie można wybrać matematyka, fizyka, informatyka albo język obcy.
    & 
    \small 
    Aby dostać się na kierunek Automatyka i Robotyka, wymagane są następujące przedmioty maturalne: Matematyka, Fizyka oraz Informatyka. Dodatkowo, w drugim przedmiocie maturalnym można wybrać Język Obcy. \\
    \bottomrule
\end{tabularx}

\begin{tcolorbox}[colback=red!5!white,colframe=red!50!black,title=Wnioski i obserwacje]
    \begin{itemize}
        \item \textbf{Utrata relacji logicznych w grafie:} GraphRAG popełnił krytyczny błąd merytoryczny. Wyciągając dane z tabeli, system zapisał przedmioty jako prostą listę encji przypisanych do kolumn P1 i P2, całkowicie gubiąc słowo "albo" (logiczne LUB/XOR). W efekcie model językowy błędnie założył, że wymagane są \textbf{wszystkie} te przedmioty naraz ("Matematyka, Fizyka oraz Informatyka").
        \item \textbf{Sukces surowego tekstu (Standard RAG):} Standardowy RAG poradził sobie z tym zapytaniem znacznie lepiej. Ponieważ silnik wektorowy zwrócił oryginalny wycinek tekstu, w którym słowo "albo" ocalało, model wygenerował w pełni poprawną odpowiedź zgodną z rzeczywistymi warunkami (kandydat wybiera jeden z tych przedmiotów).
    \end{itemize}
\end{tcolorbox}
\vspace{1em}

\newpage

\section*{Porównanie: Kierunki stacjonarne z progiem poniżej 700 punktów}

\begin{tcolorbox}[colback=blue!5!white,colframe=blue!75!black,title=Pytanie użytkownika]
    Wymień mi 3 kierunki stacjonarne, na które próg punktowy wynosił poniżej 700 punktów.
\end{tcolorbox}

\noindent
\begin{tabularx}{\textwidth}{@{} XX @{}}
    \toprule
    \textbf{Standardowy RAG} & \textbf{GraphRAG (Hybrid)} \\
    \midrule
    \small 
    Nie posiadam informacji na temat progów punktowych dla kierunków studiów. Proszę sprawdzić aktualne dane na stronie uczelni lub w dokumentach rekrutacyjnych.
    & 
    \small 
    Trzy kierunki stacjonarne, na które próg punktowy wynosił poniżej 700 punktów, to:
    \begin{enumerate}
        \item Elektrotechnika (prog\_min: 502)
        \item Inżynieria Metali Nieżelaznych (prog\_min: 341)
        \item Ceramika (prog\_min: 377)
    \end{enumerate} \\
    \bottomrule
\end{tabularx}

\begin{tcolorbox}[colback=green!5!white,colframe=green!50!black,title=Wnioski i obserwacje]
    \begin{itemize}
        \item \textbf{Przewaga danych strukturyzowanych (GraphRAG):} Model hybrydowy doskonale poradził sobie z filtrowaniem numerycznym (próg $<700$) i kategorycznym (tryb stacjonarny). Dzięki przypisaniu atrybutów do węzłów grafu (np. \texttt{prog\_min}), system potrafił zrozumieć zależność matematyczną i zwrócić poprawne wyniki.
        \item \textbf{Błąd semantyczny wyszukiwania wektorowego (Standard RAG):} Standardowy RAG zaliczył całkowitą porażkę w doborze kontekstu (Retrieval failure). Szukając słowa "punktowy/punktów", algorytm pobrał dokumenty dotyczące punktów ECTS z Regulaminu Studiów, całkowicie omijając tabelę z progami rekrutacyjnymi.
        \item \textbf{Ograniczenia wektorów w analityce:} Zwykłe bazy wektorowe mają ogromne trudności z zadaniami wymagającymi sortowania, filtrowania i operowania na ułamkach czy liczbach. Wektory sprawdzają się w szukaniu podobieństwa znaczeniowego, a nie w zapytaniach bazodanowych, co udowadnia sukces ustrukturyzowanego podejścia GraphRAG.
    \end{itemize}
\end{tcolorbox}
\vspace{1em}

\newpage

\section*{Porównanie: Warunki i kryteria stypendium socjalnego}

\begin{tcolorbox}[colback=blue!5!white,colframe=blue!75!black,title=Pytanie użytkownika]
    Kto dokładnie może ubiegać się o stypendium socjalne na AGH i jakie kryterium dochodowe / warunki formalne trzeba spełnić?
\end{tcolorbox}

\noindent
\begin{tabularx}{\textwidth}{@{} XX @{}}
    \toprule
    \textbf{Standardowy RAG} & \textbf{GraphRAG (Hybrid)} \\
    \midrule
    \small 
    O stypendium socjalne na AGH mogą ubiegać się studenci, w tym cudzoziemcy, którzy spełniają określone kryteria dochodowe oraz formalne. Kryteria dochodowe nie są dokładnie określone w kontekście, ale stypendium przysługuje w szczególnych sytuacjach, takich jak półsieroctwo, ciężka choroba członka rodziny, pobieranie zasiłku z pomocy społecznej lub zamieszkanie w domu studenckim AGH. Wymagane jest również rzetelne udokumentowanie przesłanek w dniu wydania decyzji.
    & 
    \small 
    O stypendium socjalne na AGH mogą ubiegać się studenci, w tym cudzoziemcy, na zasadach obowiązujących obywateli polskich. Kryterium dochodowe ustala się na podstawie dochodu netto na osobę w rodzinie studenta, uwzględniając dochody studenta, małżonka, dzieci oraz rodziców lub opiekunów. Dodatkowo, student może otrzymać stypendium w zwiększonej wysokości w szczególnych przypadkach, takich jak sieroctwo, ciężka choroba członka rodziny, pobieranie zasiłku z pomocy społecznej czy zamieszkanie w domu studenckim AGH. Przesłanki te muszą być udokumentowane w dniu wydania decyzji. \\
    \bottomrule
\end{tabularx}

\begin{tcolorbox}[colback=red!5!white,colframe=red!50!black,title=Wnioski i obserwacje]
    \begin{itemize}
        \item \textbf{Porażka grafu:} Silnik GraphRAG nie odnalazł żadnych węzłów dla tego zapytania (\texttt{Brak wyników}). System automatycznie przeszedł w tryb awaryjny (Fallback), opierając odpowiedź wyłącznie na tych samych słabych wektorach co standardowy RAG.
        \item \textbf{Przesunięcie uwagi:} Oba modele nie odpowiedziały poprawnie na pytanie o \textit{podstawowe} zasady stypendium socjalnego. Algorytmy wektorowe "złapały" fragmenty o stypendium w \textit{zwiększonej} wysokości oraz o cudzoziemcach ponieważ miały one silniejsze nasycenie słowami kluczowymi. W efekcie standardowy RAG zaczął wręcz twierdzić, że zwykłe stypendium przysługuje m.in. za "zamieszkanie w akademiku" (co jest absurdem i dotyczy tylko dodatku).
        \item \textbf{Ograniczenia dokumentu źródłowego:} Żaden z modeli nie podał konkretnej kwoty progu dochodowego. Wynika to z faktu, że sam Regulamin nie definiuje twardej kwoty (jest ona ogłaszana corocznym komunikatem Rektora). Modele jednak nie potrafiły tego jasno wyartykułować na bazie pobranych skrawków tekstu.
    \end{itemize}
\end{tcolorbox}
\vspace{1em}

\newpage

\section*{Porównanie: Punkty rekrutacyjne za dyplom zawodowy (Technik Informatyk)}

\begin{tcolorbox}[colback=blue!5!white,colframe=blue!75!black,title=Pytanie użytkownika]
    Posiadam dyplom technika informatyka. Czy ten egzamin zawodowy daje mi jakieś przywileje lub punkty przy rekrutacji na Informatykę i Systemy Inteligentne?
\end{tcolorbox}

\noindent
\begin{tabularx}{\textwidth}{@{} XX @{}}
    \toprule
    \textbf{Standardowy RAG} & \textbf{GraphRAG (Hybrid)} \\
    \midrule
    \small 
    Nie znalazłem informacji na temat przywilejów lub punktów przy rekrutacji na Informatykę i Systemy Inteligentne w związku z posiadanym dyplomem technika informatyka. Proszę sprawdzić regulamin rekrutacji lub skontaktować się z uczelnią w celu uzyskania dokładnych informacji.
    & 
    \small 
    Tak, dyplom technika informatyka (351203) daje punkty przy rekrutacji na kierunek Informatyka i Systemy Inteligentne jako wskaźnik P2. \\
    \bottomrule
\end{tabularx}

\begin{tcolorbox}[colback=green!5!white,colframe=green!50!black,title=Wnioski i obserwacje]
    \begin{itemize}
        \item \textbf{Precyzyjne mapowanie relacji (GraphRAG):} Model oparty na grafach wiedzy zadziałał w tym przypadku perfekcyjnie. Powiązanie konkretnego kierunku z długą listą akceptowanych kwalifikacji zawodowych to typowy problem relacyjny. Graf bezbłędnie zidentyfikował węzeł kierunku i odczytał z jego atrybutów, że dyplom "Technika informatyka" zamienia się wprost na wskaźnik rekrutacyjny P2.
        \item \textbf{Zgubny wpływ "szumu semantycznego" (Standard RAG):} Standardowy silnik wektorowy zanotował całkowitą porażkę. Szukając słów "egzamin" i "dyplom", algorytm pobrał z bazy fragmenty \textit{Regulaminu Studiów} opowiadające o egzaminach dyplomowych na koniec nauki, zamiast sięgnąć do zasad rekrutacji.
        \item \textbf{Ograniczenia wektorów przy listach:} Ten przykład świetnie udowadnia, że RAG wektorowy słabo radzi sobie z długimi wykazami (np. tabelami zawodów uprawniających do przyjęcia na dany kierunek). Jeśli słowo kluczowe nie leży fizycznie blisko nazwy kierunku w tekście źródłowym, wektor tego nie połączy, podczas gdy graf tworzy bezpośrednie i nierozerwalne "mosty" między takimi informacjami.
    \end{itemize}
\end{tcolorbox}
\vspace{1em}

\newpage

\section*{Porównanie: Obliczanie wskaźnika rekrutacyjnego (WR)}

\begin{tcolorbox}[colback=blue!5!white,colframe=blue!75!black,title=Pytanie użytkownika]
    Jak obliczyć wskaźnik rekrutacyjny dla kandydatów z nową maturą?
\end{tcolorbox}

\noindent
\begin{tabularx}{\textwidth}{@{} XX @{}}
    \toprule
    \textbf{Standardowy RAG} & \textbf{GraphRAG (Hybrid)} \\
    \midrule
    \small 
    Aby obliczyć wskaźnik rekrutacyjny WR dla kandydatów z nową maturą, należy zastosować wzór:
    
    $$WR = 2 \cdot M + 6 \cdot P1 + 2 \cdot P2 + D$$
    
    gdzie:
    \begin{itemize}
        \item M to liczba punktów uzyskanych z przeliczenia wyników z matematyki,
        \item P1 i P2 to liczby punktów uzyskanych z przeliczenia wyników z przedmiotów uwzględnionych w tabeli 1,
        \item D to dodatkowe punkty, jeśli są przyznawane.
    \end{itemize}
    
    Należy pamiętać, że do P1 i P2 wybierane są wyniki najkorzystniejsze, a przedmioty nie mogą być takie same.
    & 
    \small 
    Aby obliczyć wskaźnik rekrutacyjny (WR) dla kandydatów z nową maturą, należy skorzystać z następującego wzoru:

    $$WR = 2 * M + 6 * P1 + 2 * P2 + D$$

    gdzie:
    \begin{itemize}
        \item M – liczba punktów procentowych uzyskanych w części pisemnej egzaminu maturalnego z matematyki na poziomie podstawowym,
        \item P1, P2 – liczby punktów procentowych uzyskanych z części pisemnej egzaminu maturalnego na poziomie rozszerzonym z przedmiotu wymienionego w Tabeli 1,
        \item D – uwzględniane osiągnięcia sportowe lub dodatkowe osiągnięcia kandydata.
    \end{itemize}
    
    Maksymalna wartość WR wynosi 1000 punktów. \\
    \bottomrule
\end{tabularx}

\begin{tcolorbox}[colback=green!5!white,colframe=green!50!black,title=Wnioski i obserwacje]
    \begin{itemize}
        \item \textbf{Skuteczność obu metod:} Oba systemy poprawnie zidentyfikowały i odtworzyły prawidłowy wzór na wskaźnik rekrutacji (WR). Jest to spodziewane, ponieważ główny wzór występuje centralnie w dokumentach i łatwo go namierzyć zarówno wektorowo, jak i przez encję grafową ("wzór").
        \item \textbf{Precyzja zmiennych w GraphRAG:} GraphRAG dostarczył znacznie bardziej precyzyjne definicje zmiennych (np. M -- poziom \textit{podstawowy}, P1/P2 -- poziom \textit{rozszerzony}). Standardowy RAG "zgubił" informację o poziomach matury z powodu faktu, że skleił ze sobą definicje dla matury międzynarodowej (EB/IB) widoczne w swoich logach źródłowych.
        \item \textbf{Dopełnienie kontekstu (Standard RAG):} Standardowy RAG podał przydatną zasadę o "najkorzystniejszym wyniku" i "nie powtarzaniu przedmiotów", która znajdowała się w innym paragrafie wokół pobranego wektora. Z kolei GraphRAG podał maksymalną pulę punktów (1000 pkt), co wynikało ze zwięzłego opisu właściwości samego węzła.
    \end{itemize}
\end{tcolorbox}
\vspace{1em}

\newpage

\section*{Porównanie: Obliczenie szans na przyjęcie na Informatykę Stosowaną}

\begin{tcolorbox}[colback=blue!5!white,colframe=blue!75!black,title=Pytanie użytkownika]
    Na bazie wzorów rekrutacyjnych i progów dla Informatyki stosowanej, oblicz, czy dostanę się na studia, jeśli z matematyki rozszerzonej mam 79\%, a z informatyki rozszerzonej 65\%, a z języka angielskiego rozszerzonego 80\%, a z matematyki na podstawowym 100\%?
\end{tcolorbox}

\noindent
\begin{tabularx}{\textwidth}{@{} XX @{}}
    \toprule
    \textbf{Standardowy RAG} & \textbf{GraphRAG (Hybrid)} \\
    \midrule
    \small 
    Nie posiadam wystarczających informacji na temat progów punktowych dla Informatyki stosowanej, aby obliczyć, czy się dostaniesz. Proszę sprawdzić aktualne progi na stronie rekrutacji uczelni.
    & 
    \small 
    Aby obliczyć wskaźnik rekrutacji (WR) dla kierunku Informatyka Stosowana, użyjemy wzoru:
    
    $$WR = 2 * M + 6 * P1 + 2 * P2 + D$$
    
    Gdzie:
    \begin{itemize}
        \item $M$ to wynik z matematyki na poziomie podstawowym (100\%),
        \item $P1$ to wynik z matematyki rozszerzonej (79\%),
        \item $P2$ to wynik z informatyki rozszerzonej (65\%),
        \item $D$ to dodatkowe osiągnięcia (zakładamy, że $D = 0$, jeśli nie ma dodatkowych osiągnięć).
    \end{itemize}
    
    Podstawiając wartości:
    $$WR = 2 * 100 + 6 * 79 + 2 * 65 + 0$$
    $$WR = 200 + 474 + 130$$
    $$WR = 804$$
    
    Minimalny próg punktowy dla Informatyki Stosowanej wynosi 829 punktów. Zatem, z wynikiem 804 punktów nie dostaniesz się na studia. \\
    \bottomrule
\end{tabularx}

\begin{tcolorbox}[colback=green!5!white,colframe=green!50!black,title=Wnioski i obserwacje]
    \begin{itemize}
        \item \textbf{Precyzja i wieloetapowe rozumowanie (Multi-hop reasoning):} GraphRAG poradził sobie ze znalezieniem abstrakcyjnego wzoru i poprawnie zmapował wyniki użytkownika do zmiennych matematycznych, wykonując obliczenia.
        \item \textbf{Ograniczenia wyszukiwania wektorowego:} Standardowy RAG oparty tylko na podobieństwie tekstu odnalazł dokumenty z progami dla innych kierunków (np. Informatyka Medyczna, Informatyka Techniczna). Ponieważ nie znalazł "Informatyki Stosowanej" w topowych fragmentach, zrezygnował z odpowiedzi, ignorując nawet możliwość podania samego obliczonego wskaźnika.
    \end{itemize}
\end{tcolorbox}
\vspace{1em}

\newpage

\section*{Porównanie: Przeniesienie na inny kierunek i statystyki rekrutacyjne}

\begin{tcolorbox}[colback=blue!5!white,colframe=blue!75!black,title=Pytanie użytkownika]
    Jeśli dostałem się na Budownictwo stacjonarne, ale po 1 semestrze uznam, że to błąd i chcę przenieść się na Inżynierię Materiałową (ile na tym kierunku było osób na jedno miejsce?), to czy Regulamin Studiów AGH pozwala na takie przeniesienie i co muszę zrobić?
\end{tcolorbox}

\noindent
\begin{tabularx}{\textwidth}{@{} XX @{}}
    \toprule
    \textbf{Standardowy RAG} & \textbf{GraphRAG (Hybrid)} \\
    \midrule
    \small 
    Możesz przenieść się na Inżynierię Materiałową po uzyskaniu wpisu na co najmniej drugi semestr. Musisz złożyć wniosek do Dziekana Wydziału przyjmującego oraz spełnić wszystkie obowiązki wobec dotychczasowego Wydziału. Informacja o liczbie osób na jedno miejsce na kierunku Inżynieria Materiałowa nie została podana w dostępnych materiałach.
    & 
    \small 
    Regulamin Studiów AGH pozwala na przeniesienie się na inny kierunek, w tym z Budownictwa na Inżynierię Materiałową, pod warunkiem uzyskania zgody Dziekana Wydziału przyjmującego. Musisz złożyć wniosek o przeniesienie po ukończeniu co najmniej jednego semestru, a także spełnić wszystkie obowiązki wobec uczelni i zaliczyć dotychczasowe moduły zajęć. 
    
    Na kierunku Inżynieria Materiałowa było 176 miejsc na 60 kandydatów, co oznacza, że na jedno miejsce przypada około 2,93 osoby. \\
    \bottomrule
\end{tabularx}

\begin{tcolorbox}[colback=green!5!white,colframe=green!50!black,title=Wnioski i obserwacje]
    \begin{itemize}
        \item \textbf{Obsługa zapytań wielowątkowych (Multi-intent retrieval):} GraphRAG poradził sobie rewelacyjnie ze znalezieniem odpowiedzi z dwóch zupełnie różnych domen (regulamin studiów oraz tabela z progami punktowymi). Graf zidentyfikował węzeł z kierunkiem i wyciągnął atrybut liczby miejsc, jednocześnie zachowując kontekst prawny z wektorów.
        \item \textbf{Gubienie wątków przez Standardowy RAG:} Zwykły system oparty na wektorach skupił się na dominującej semantycznie części zapytania (Regulamin, przeniesienie, semestr) i odnalazł świetne fragmenty regulaminu. Całkowicie jednak zignorował wtrącenie w nawiasie dotyczące statystyk, uznając je za szum (tzw. "Missed intent").
        \item \textbf{Ciekawy błąd generowania języka (NLG Slip) w GraphRAG:} Mimo że GraphRAG poprawnie policzył stosunek $176 / 60 \approx 2,93$, w zdaniu wygenerowanym przez LLM pomylono pojęcia, pisząc "176 miejsc na 60 kandydatów" zamiast "176 kandydatów na 60 miejsc". To przykład na to, że nawet jeśli ustrukturyzowane dane wyciągnięte z grafu (miejsca: 176/60) są poprawne, model językowy formujący ostateczną wypowiedź może popełnić nielogiczny błąd semantyczny.
    \end{itemize}
\end{tcolorbox}
\vspace{1em}, dodaj mi pliki tex ustruktyryzowane na modłę tego, nie zmieniaj wiecej nic, w folderze RAG_benchmark, bazuj to na podstawie tego folderu: evaluation\results\compare\tracesC:\Users\wolny\Documents\projects\agh\6sem\lrag-assistant\evaluation\results\compare\traces